% Source printed pp. 476-489 / PDF pp. 484-497.
\section{几何性质}\label{src:9.4}
\begin{epigraph}
有人看见它, 这是奇迹; 另有人谈论它, 这是奇迹; 又有人听见它, 这也是奇迹. 然而即使听过, 也无人真正知晓它.
\epigraphauthor{``薄伽梵歌''}
\end{epigraph}
我们证明半稳定向量丛模叠 $\cBun_{r,d}^{\mathrm{ss}}$, $\cBun_{r,\xi}^{\mathrm{ss}}$ 及其模空间 $M_{r,d}^{\mathrm{ps}}$, $M_{r,\xi}^{\mathrm{ps}}$ 的几何性质: 不可约性 (定理~\bookref{9.4.1}), 单有理性 (定理~\bookref{9.4.2}), 局部阶乘性 (命题~\bookref{9.4.12}), 行列式线丛的性质 (\S\bookref{9.4.3-subsection}), Poincar\'e 丛的刻画 (定理~\bookref{9.4.34}), Picard 群的计算 (定理~\bookref{9.4.36}) 和上同调 (\S\bookref{9.4.6-subsection}). 除 $M_{r,\xi}^{\mathrm{ps}}$ 为 Fano 簇 (推论~\bookref{9.4.43}) 外, 本节的结果与证明不依赖于 $M_{r,d}^{\mathrm{ss}}$ 或 $M_{r,\xi}^{\mathrm{ss}}$ 的射影性.
\subsection{半稳定模叠的不可约性与固定行列式模叠的单有理性}\label{src:9.4.1-subsection}
\begin{theorem}[\texorpdfstring{$\cBun_{r,d}^{\mathrm{ss}}$}{Bunss} 的不可约性]\label{src:9.4.1}
设 $C$ 是代数闭域 $\Bbbk$ 上亏格 $g\geq2$ 的光滑连通射影曲线. 对整数 $r>0$ 和 $d$, $\cBun_{r,d}^{\mathrm{ss}}$ 不可约. 特别地, $M_{r,d}^{\mathrm{ps}}$ 不可约.
\end{theorem}
\begin{theorem}[\texorpdfstring{$\cBun_{r,\xi}^{\mathrm{ss}}$}{Bunss} 的单有理性]\label{src:9.4.2}
设 $C$ 是代数闭域 $\Bbbk$ 上亏格 $g\geq2$ 的光滑连通射影曲线. 对 $r>0$ 和 $\xi\in\mathrm{Pic}(C)$, $\cBun_{r,\xi}^{\mathrm{ss}}$ 不可约且单有理, 即存在支配有理映射 $\mathbb P^N\dashrightarrow\cBun_{r,\xi}^{\mathrm{ss}}$. 特别地, $M_{r,\xi}^{\mathrm{ps}}$ 不可约且单有理.
\end{theorem}
两个定理的关键是下述引理, 它将半稳定向量丛表为其行列式被平凡丛的扩张. 它也是计算 $\operatorname{Pic}(\cBun_{r,d}^{\mathrm{ss}})$ (定理~\bookref{9.4.36}) 的关键. 论证基于 Serre 的思想 (见 \cite[\S2]{Ati57}).
\begin{lemma}\label{src:9.4.3}
$C$ 上每个秩 $r>0$, 次数 $d>r(2g-1)$ 的半稳定向量丛 $F$ 都可写为扩张
\[
0\to\mathcal O_C^{\oplus r-1}\to F\to\det F\to0.
\]
\end{lemma}
\begin{proof}
因为 $\mu(F)>2g-1$, $F$ 由整体截面生成 (引理~\bookref{9.2.7}). 断言对一般的维数 $r-1$ 子空间 $W\subseteq\mathrm H^0(C,F)$, 求值映射
\[
\operatorname{ev}_W:W\otimes\mathcal O_C\hookrightarrow\mathrm H^0(C,F)\otimes\mathcal O_C\twoheadrightarrow F
\]
的余核为线丛. 因而 $\operatorname{ev}_W$ 单射, 检查所诱导的短正合列的行列式可知其余核必为 $\det F$. 用维数计数建立此断言.

可将每个子空间 $W$ 看作 Grassmann 簇 $\Gr(k,\mathrm H^0(C,F))$ 的元素, 其中 $k=\mathrm h^0(C,F)-(r-1)$; 回顾我们约定 Grassmann 簇参数化商 $\mathrm H^0(C,F)\twoheadrightarrow\mathrm H^0(C,F)/W$. 我们将证明
\begin{equation}\label{src:9.4.4}
\{W\mid\operatorname{coker}(\operatorname{ev}_W)\text{ 不是线丛}\}\subseteq\Gr(k,\mathrm H^0(C,F))
\end{equation}
的余维为 $1$ (为什么此轨迹闭?). 求值映射在点 $p\in C(\Bbbk)$ 处的纤维 $\operatorname{ev}_{W,p}:W\to F\otimes\kappa(p)$ 单射, 当且仅当 $\operatorname{coker}(\operatorname{ev}_W)$ 在 $p$ 的某邻域中为线丛. 轨迹
\begin{equation}\label{src:9.4.5}
\{W\mid e_{W,p}\text{ 单射}\}\subseteq\Gr(k,\mathrm H^0(C,F))
\end{equation}
的余维为 $2$ (你知道为什么吗?). 让 $p\in C$ 变化, 推出 \bookeqref{9.4.4} 的余维为 $1$. 另见 \cite[引理 7.1]{DN89} 和 \cite[引理 8.5.3]{LP97}.
\end{proof}
\begin{exercise}[中等]\label{src:9.4.6}
验证省略的细节: 轨迹 \bookeqref{9.4.4} 闭, 轨迹 \bookeqref{9.4.5} 的余维为 $2$.
\end{exercise}
\begin{exercise}[困难, 后文需要]\label{src:9.4.7}
证明 \bookeqref{9.4.4} 是次数为 $d$ 的超曲面.

\textit{提示:} 应用 Porteous 公式 (见 \cite[\S14.4]{Ful98} 或 \cite[\S12.2]{EH16}).
\end{exercise}
\begin{proof}[单有理性 (\bookref{9.4.2}) 的证明]
对 $p\in C(\Bbbk)$, 与 $\mathcal O_C(p)$ 作张量积定义同构 $\cBun_{r,\xi}^{\mathrm{ss}}\to\cBun_{r,\xi(rp)}^{\mathrm{ss}}$. 反复应用此同构后, 可假设 $\deg\xi$ 任意大, 例如大于 $r(2g-1)$. 由习题~\bookref{1.3.22}, 仿射空间 $\underline{\operatorname{Ext}}_{\mathcal O_C}^1(\xi,\mathcal O_C^{\oplus r-1}):=\mathbb A(\operatorname{Ext}_{\mathcal O_C}^1(\xi,\mathcal O_C^{\oplus r-1})^\vee)$ 表示扩张函子. $C\times\underline{\operatorname{Ext}}_{\mathcal O_C}^1(\xi,\mathcal O_C^{\oplus r-1})$ 上的万有扩张 $\mathcal E_{\mathrm{univ}}$ 定义映射
\[
\underline{\operatorname{Ext}}_{\mathcal O_C}^1(\xi,\mathcal O_C^{\oplus r-1})\to\cBun_{r,\xi}.
\]
由引理~\bookref{9.4.3}, 此映射的像包含 $\cBun_{r,\xi}^{\mathrm{ss}}$.
\end{proof}
\begin{remark}[有理性]\label{src:9.4.8}
当 $\gcd(r,d)=1$ 时, 模空间 $M_{r,d}^{\mathrm{ps}}$ 有理; 见 \cite{New75,New80,KS99,Hof10}. 一般情形下有理性仍是未解决的问题.
\end{remark}
\begin{proof}[不可约性 (\bookref{9.4.1}) 的证明]
对点 $p\in C(\Bbbk)$ 反复应用同构 $\cBun_{r,d}^{\mathrm{ss}}\to\cBun_{r,dr}^{\mathrm{ss}}$, $E\mapsto E(p)$, 可假设 $d>r(2g-1)$. 由 $F\mapsto\det F$ 定义的到 Picard 叠的态射 $\cBun_{r,d}^{\mathrm{ss}}\to\underline{\mathcal P\!ic}^d$ 满射, 在 $\xi\in\underline{\mathcal P\!ic}^d(\Bbbk)$ 上的纤维为 $\operatorname{Ext}_{\mathcal O_C}^1(\xi,\mathcal O_C^{\oplus r-1})$ 的开子空间. 因为 $\underline{\mathcal P\!ic}^d$ 和每个纤维都不可约, $\cBun_{r,d}^{\mathrm{ss}}$ 也不可约.
\end{proof}
\begin{exercise}\label{src:9.4.9}
利用单有理性证明中的构造的相对版本, 给出 $\cBun_{r,d}^{\mathrm{ss}}$ 的不可约性的另一种证明. 具体地, 考虑投影 $p_2:C\times\underline{\mathcal P\!ic}^d\to\underline{\mathcal P\!ic}^d$ 及 $C\times\underline{\mathcal P\!ic}^d$ 上的万有线丛 $\xi_{\mathrm{univ}}$, 定义
\[
\mathbb E:=\left(p_{2,*}\mathcal E\!xt_{\mathcal O_{C\times\underline{\mathcal P\!ic}^d}}^1
(\xi_{\mathrm{univ}},\mathcal O_{C\times\underline{\mathcal P\!ic}^d}^{\oplus r-1})\right)^\vee,
\]
并证明万有扩张 (见习题~\bookref{9.2.12}) 定义态射 $\mathbb A(\mathbb E)\to\cBun_{r,d}$, 其像包含半稳定轨迹. 另见 \cite[定理 8.5.2]{LP97}.
\end{exercise}
\subsection{模空间的奇点}\label{src:9.4.2-subsection}
半稳定向量丛模叠 $\cBun_{r,d}^{\mathrm{ss}}$ 和 $\cBun_{r,\xi}^{\mathrm{ss}}$ 光滑. 由良模空间的奇点的一般结果 (\S\bookref{7.5.11-subsection}), 模空间 $M_{r,d}^{\mathrm{ps}}$ 和 $M_{r,\xi}^{\mathrm{ps}}$ 正规, Cohen--Macaulay, 且有有理奇点. 在稳定轨迹上, $\cBun_{r,d}^{\mathrm s}\to M_{r,d}^{\mathrm s}$ 和 $\cBun_{r,\xi}^{\mathrm s}\to M_{r,\xi}^{\mathrm s}$ 分别为带 $\mathbb G_m$-带和 $\boldsymbol\mu_r$-带的 gerbe, 因而 $M_{r,d}^{\mathrm s}$ 与 $M_{r,\xi}^{\mathrm s}$ 也光滑. 特别地, 若 $\gcd(r,d)=1$, 则 $M_{r,d}^{\mathrm{ps}}$ 和 $M_{r,\xi}^{\mathrm{ps}}$ 光滑.

由局部结构定理 (\bookref{7.7.1}) 以及 $\cBun_{r,d}$ 在向量丛 $F$ 处的切空间与 $\operatorname{Ext}_{\mathcal O_C}^1(F,F)$ 的等同, 模叠 $\cBun_{r,d}$ 在 $F$ 处平展局部同构于
\[
[\underline{\operatorname{Ext}}_{\mathcal O_C}^1(F,F)/\underline{\operatorname{Aut}}(F)],
\]
其中 $\underline{\operatorname{Ext}}_{\mathcal O_C}^1(F,F):=\mathbb A(\operatorname{Ext}_{\mathcal O_C}^1(F,F)^\vee)$, 作用为
\[
\alpha\cdot[0\to F\xrightarrow aE\xrightarrow bF\to0]
=[0\to F\xrightarrow{a\circ\alpha^{-1}}E\xrightarrow{\alpha\circ b}F\to0].
\]
若 $F=\bigoplus F_i^{\oplus r_i}$ 多重稳定, $F_i$ 是互不同构的稳定向量丛, 则 $\operatorname{Aut}(F)\cong\prod_i\mathrm{GL}_{r_i}$ 在向量空间
\[
\operatorname{Ext}_{\mathcal O_C}^1(F,F)\cong
\bigoplus_{i,j}\operatorname{Ext}_{\mathcal O_C}^1(F_i,F_j)\otimes\mathrm{Mat}_{r_i,r_j}
\]
上的作用为 $(g_i)\cdot(e_{ij}\otimes M_{ij})=e_{ij}\otimes g_iM_{ij}g_j^{-1}$. 此外, 由良模空间的局部结构定理 (\bookref{7.7.3}), 模空间 $M_{r,d}^{\mathrm{ps}}$ 在多重稳定向量丛 $F$ 处平展局部同构于 $\underline{\operatorname{Ext}}_{\mathcal O_C}^1(F,F)\sslash\underline{\operatorname{Aut}}(F)$. 利用习题~\bookref{9.2.30} 给出的切空间等同, 可类似地描述 $M_{r,\xi}^{\mathrm{ps}}$ 的局部结构.
\begin{exercise}\label{src:9.4.10}
设 $r\geq2$, $\xi$ 为 $C$ 上次数 $d$ 的线丛.
\begin{enumerate}[(a)]
\item 若 $r=g=2$ 且 $d$ 为偶数, 证明 $M_{r,d}^{\mathrm{ps}}$ 和 $M_{r,\xi}^{\mathrm{ps}}$ 光滑.

\textit{提示:} 证明每个局部 GIT 商都与仿射空间等同; 当 $F\cong L\oplus L$ 时, 用例~\bookref{7.5.9} 确定不变量环.
\item 假设 $r\ne2$, $g\ne2$ 或 $d$ 为奇数. 证明 $M_{r,d}^{\mathrm{ps}}$ 和 $M_{r,\xi}^{\mathrm{ps}}$ 的奇点轨迹恰为严格半稳定轨迹的像.
\end{enumerate}
\end{exercise}
\begin{remark}\label{src:9.4.11}
除了上述性质, Narasimhan 与 Ramanan 在 \cite{NR69} 中还对 $g=2$ 时的 $M_{2,\xi}^{\mathrm{ps}}$ 给出显式描述: 若 $d$ 偶, $M_{2,d}^{\mathrm{ps}}\cong\mathbb P^3$; 若 $d$ 奇, $M_{2,d}^{\mathrm{ps}}$ 是 $\mathbb P^5$ 中两个二次超曲面的交. 对任意亏格, \cite{Las96} 描述 $M_{2,\xi}^{\mathrm{ps}}$ 的奇点 (例如重数与切锥).
\end{remark}
\begin{proposition}\label{src:9.4.12}
设 $r>0$ 和 $d$ 为整数, $\xi$ 为 $C$ 上线丛. 模空间 $M_{r,d}^{\mathrm{ps}}$ 和 $M_{r,\xi}^{\mathrm{ps}}$ 局部阶乘.
\end{proposition}
\begin{proof}
$r=g=2$ 时可显式验证 (见习题~\bookref{9.4.10}). 否则, 由习题~\bookref{9.2.18}, 已知 $\operatorname{codim}(\cBun_{r,d}^{\mathrm{ss}}\setminus\cBun_{r,d}^{\mathrm s})\geq2$. 考虑 $\mathbb G_m$-刚性化, 有 Cartesian 图
\[
\begin{tikzcd}[ampersand replacement=\&]
\cBun_{r,d}^{\mathrm s}\!\sslash\mathbb G_m\arrow[r]\arrow[d,"\sim"']\&\cBun_{r,d}^{\mathrm{ss}}\!\sslash\mathbb G_m\arrow[d]\\
M_{r,d}^{\mathrm s}\arrow[r]\&M_{r,d}^{\mathrm{ss}},
\end{tikzcd}
\]
其中左侧竖直映射为同构, 右侧竖直映射为良模空间. 由习题~\bookref{7.5.39}, $M_{r,d}^{\mathrm{ss}}$ 阶乘, 当且仅当 $\cBun_{r,d}^{\mathrm{ss}}\!\sslash\mathbb G_m$ 上每个线丛都下降到 $M_{r,d}^{\mathrm{ss}}$. 后一陈述正是下述引理~\bookref{9.4.13} 的结论. 另见 \cite[定理 A]{DN89}.
\end{proof}
\begin{lemma}\label{src:9.4.13}
设 $L$ 是 $\cBun_{r,d}^{\mathrm{ss}}$ 上线丛. 对每个秩 $r$, 次数 $d$ 的多重稳定向量丛 $F$, $\underline{\operatorname{Aut}}(F)/\mathbb G_m$ 在纤维 $L\otimes\kappa([F])$ 上平凡作用. 特别地, $\operatorname{Pic}(M_{r,d}^{\mathrm{ps}})\cong\operatorname{Pic}(\cBun_{r,d}^{\mathrm{ss}}\!\sslash\mathbb G_m)$.
\end{lemma}
\begin{proof}
令 $k=\gcd(r,d)$, $r'=r/k$, $d'=d/k$. 因为每个 $F_i$ 的斜率为 $d/r$, 可写为 $r_i:=\operatorname{rk}F_i=k_ir'$ 和 $d_i:=\deg F_i=k_id'$, 其中 $k_i$ 为整数. 由于 $(r,d)=\sum_i n_i(\operatorname{rk}F_i,\deg F_i)=\sum_i n_ik_i(r',d')$, 有 $k=\sum_i n_ik_i$. 利用稳定丛的存在性 (\bookref{9.2.16}) 选取稳定丛 $G\in\cBun_{r',d'}^{\mathrm s}$. 则 $G^{\oplus k}$ 是秩为 $r$, 次数为 $d$ 的多重稳定向量丛, 自同构群为 $\mathrm{GL}_k$. 因为 $\mathrm{GL}_k/\mathbb G_m=\mathrm{PGL}_k$ 没有特征标, 它在 $L\otimes\kappa([G^{\oplus k}])$ 上的作用平凡. 为说明对 $F$ 也成立, 考虑态射
\begin{equation}\label{src:9.4.14}
\begin{aligned}
\cBun_{r_1,d_1}^{\mathrm{ss}}\times\cdots\times\cBun_{r_l,d_l}^{\mathrm{ss}}&\to\cBun_{r,d}^{\mathrm{ss}}\\
(E_1,\ldots,E_l)&\mapsto E_1^{\oplus n_1}\oplus\cdots\oplus E_l^{\oplus n_l}.
\end{aligned}
\end{equation}
注意 $F$ 是 $(F_1,\ldots,F_l)$ 的像, $G^{\oplus k}$ 是 $(G^{\oplus k_1},\ldots,G^{\oplus k_l})$ 的像. 有自同构群的包含关系
\[
\operatorname{Aut}(F)\cong\prod_i\mathrm{GL}_{n_i}\subseteq
\prod_i\mathrm{GL}_{n_ik_i}\subseteq\mathrm{GL}_k=\operatorname{Aut}(G^{\oplus k}).
\]
对 \bookeqref{9.4.14} 的像中的每个向量丛 $E=E_1^{\oplus n_1}\oplus\cdots\oplus E_l^{\oplus n_l}$, 有 $\underline{\operatorname{Aut}}(F)/\mathbb G_m\subseteq\operatorname{Aut}(E)/\mathbb G_m$ 在 $L\otimes\kappa([E])$ 上的作用, 从而诱导 $\underline{\operatorname{Aut}}(F)/\mathbb G_m$ 的特征标. 因为特征标离散, 而 $\cBun_{r_1,d_1}^{\mathrm{ss}}\times\cdots\times\cBun_{r_l,d_l}^{\mathrm{ss}}$ 连通, 此特征标必与 $E$ 无关. 取 $E=G^{\oplus k}$, 因为 $\underline{\operatorname{Aut}}(F)/\mathbb G_m$ 在 $L\otimes\kappa([G^{\oplus k}])$ 上的作用平凡, 它在 $L\otimes\kappa([F])$ 上的作用也平凡. 将命题~\bookref{7.5.36} 应用于良模空间 $\cBun_{r,d}^{\mathrm{ss}}\!\sslash\mathbb G_m\to M_{r,d}^{\mathrm{ps}}$, 得到 $\cBun_{r,d}^{\mathrm{ss}}\!\sslash\mathbb G_m$ 上每个线丛都下降为 $M_{r,d}^{\mathrm{ps}}$ 上线丛, 因而 $\operatorname{Pic}(M_{r,d}^{\mathrm{ps}})\cong\operatorname{Pic}(\cBun_{r,d}^{\mathrm{ss}}\!\sslash\mathbb G_m)$.
\end{proof}
